\documentclass{article}

% ── Encoding & fonts ─────────────────────────────────────────────────────────
\usepackage[utf8]{inputenc}
\usepackage[T1]{fontenc}
\usepackage{lmodern}

% ── Page geometry ─────────────────────────────────────────────────────────────
\usepackage[a4paper, margin=2.5cm]{geometry}

% ── Hyperlinks ────────────────────────────────────────────────────────────────
\usepackage[hidelinks, unicode]{hyperref}

% ── Source code listings ─────────────────────────────────────────────────────
\usepackage{listings}
\lstset{
  basicstyle=\ttfamily\small,
  breaklines=true,
  keepspaces=true,
  columns=flexible,
  frame=single,
  numbers=left,
  numberstyle=\tiny,
}

% ── List customisation ───────────────────────────────────────────────────────
\usepackage{enumitem}

% ── Graphics ─────────────────────────────────────────────────────────────────
\usepackage{graphicx}

% ── Tables ───────────────────────────────────────────────────────────────────
\usepackage{booktabs}
\usepackage{tabularx}

% ── Document metadata ────────────────────────────────────────────────────────
\title{Advances in Markdown Processing\\[0.5em]{\normalfont\large A Comprehensive Survey}}
\author{Andreas W. Weber}
\date{2026-05-03}

\begin{document}
\maketitle

\tableofcontents

\listoffigures

\listoftables

\lstlistoflistings

\section{Introduction}

Markdown-to-LaTeX conversion has been studied extensively~\cite{knuth1986texbook}. The \texttt{TeX} typesetting system~\cite{lamport1994latex} remains the standard for academic publishing, while Markdown~\cite{gruber2004markdown} has become the dominant lightweight markup language for technical writing. This paper presents MarkTeX, a transpiler that bridges the two.

The CommonMark specification~\cite{commonmark2021spec} defines the authoritative Markdown grammar; our parser targets a strict subset of it, extended with custom syntax for cross-references and document metadata.

\subsection{Motivation}

Three problems motivate this work:

\begin{enumerate}
  \item Existing tools~\cite{knuth1986texbook} require \textbf{complex configuration} to produce well-structured LaTeX output.
  \item Reference management is \textit{manual} in most pipelines.
  \item \textbf{\textit{No open standard}} exists for embedding LaTeX-specific constructs in Markdown source.
\end{enumerate}

The following sections address each problem in turn.


\noindent\rule{\linewidth}{0.4pt}

\section{Architecture}

As shown in Figure~\ref{fig:fig_pipeline}, the MarkTeX pipeline consists of three stages: \textbf{lexical analysis}, \textit{AST construction}, and \textbf{\textit{code generation}}.

\begin{figure}[h]
  \centering
  \includegraphics[width=0.9\linewidth]{figures/pipeline.png}
  \caption{The three-stage transpilation pipeline, adapted from~\cite{gruber2004markdown}.}
  \label{fig:fig\_pipeline}
\end{figure}

\subsection{Parser}

The parser follows the two-phase approach of~\cite{commonmark2021spec}: a block pass identifies structural elements, and an inline pass handles emphasis, code spans, and custom references. The grammar is implemented without backtracking; ambiguous constructs default to the CommonMark resolution rules.

Block-level elements recognised by the parser:

\begin{itemize}
  \item \textbf{Headings}: ATX (\texttt{\# H1} through \texttt{\#\#\#\#\#\# H6}) and setext (\texttt{===}, \texttt{---})
  \item \textbf{Lists}: bullet (\texttt{-}, \texttt{*}, \texttt{+}) and ordered (\texttt{1.}, \texttt{2.}, …) with tight/loose detection
  \item \textbf{Code}: fenced (\texttt{```}) and indented (4-space)
  \item \textbf{Tables}: GFM pipe tables with \texttt{:---}, \texttt{---:}, \texttt{:---:} alignment
  \item \textbf{Block quotes} and \textbf{thematic breaks}
\end{itemize}

Inline elements recognised:

\begin{itemize}
  \item Emphasis (\texttt{*}, \texttt{\_}), strong (\texttt{**}), strong-emphasis (\texttt{***})
  \item Code spans (\texttt{`}), strikethrough (\texttt{\textasciitilde{}\textasciitilde{}})
  \item Links \texttt{[text](url)} and images \texttt{![alt](url)}
  \item Custom refs: \texttt{[C\#key]}, \texttt{[F\#key]}, \texttt{[T\#key]}
\end{itemize}

\subsection{AST}

The internal AST is depicted in Figure~\ref{fig:fig_ast}. Every node carries a source position for error reporting.

\begin{figure}[ht]
  \centering
  \includegraphics[width=0.7\linewidth]{figures/ast.png}
  \caption{AST node hierarchy. Block nodes are shown above the dashed line; inline nodes below. See Table~\ref{tab:tab_comparison} for node-count statistics.}
  \label{fig:fig\_ast}
\end{figure}

Each node type implements a \texttt{Type() NodeType} method. Traversal is handled by a single \texttt{Walk} function; visitors implement \texttt{Enter*/Leave*} pairs and embed \texttt{BaseVisitor} for default no-op behaviour.

\subsubsection{Extension Mechanism}

\begin{quote}
Adding a new construct requires five deterministic steps: define the AST node, extend the \texttt{Visitor} interface, provide a \texttt{BaseVisitor} no-op, register the dispatch case in \texttt{Walk}, and implement the generator method. The pattern was validated against~\cite{lamport1994latex} and~\cite{commonmark2021spec}.

\end{quote}


\noindent\rule{\linewidth}{0.4pt}

\section{Evaluation}

\subsection{Feature Comparison}

Table~\ref{tab:tab_comparison} compares MarkTeX against existing tools.

\begin{table}[ht]
\centering
\caption{Feature comparison of Markdown-to-LaTeX tools. Based on criteria from~\cite{knuth1986texbook} and~\cite{lamport1994latex}.}
\begin{tabular}{lcccr}
\toprule
Tool & Typed AST & Extensions & Cross-refs & Speed \\
\midrule
MarkTeX & yes & yes & yes & 12 ms \\
pandoc & yes & yes & partial & 85 ms \\
markdown & no & no & no & 3 ms \\
kramdown & partial & limited & no & 22 ms \\
\bottomrule
\end{tabular}
\label{tab:tab\_comparison}
\end{table}

\subsection{Performance}

Table~\ref{tab:tab_performance} reports processing times on a corpus of 1 000 documents.

\begin{table}[h]
\centering
\caption{Processing time by document size (adapted from~\cite{knuth1986texbook}). See Figure~\ref{fig:fig_pipeline} for pipeline details.}
\begin{tabular}{cccr}
\toprule
Input size & Parse time & Generate time & Total \\
\midrule
1 KB & 0.5 ms & 0.2 ms & 0.7 ms \\
10 KB & 4.2 ms & 1.8 ms & 6.0 ms \\
100 KB & 38 ms & 15 ms & 53 ms \\
\bottomrule
\end{tabular}
\label{tab:tab\_performance}
\end{table}


\noindent\rule{\linewidth}{0.4pt}

\section{Implementation}

\subsection{Code Example}

The public API consists of three functions:

\begin{lstlisting}[language=go]
// Parse Markdown with all extensions enabled.
doc := parser.Parse(src, parser.ExtAll)

// Generate LaTeX fragment output.
err := generator.Generate(doc, generator.Options{}, os.Stdout)

// Or use the high-level transpiler.
out, err := transpiler.TranspileString(src, transpiler.Options{
    Standalone: true,
    Extensions: transpiler.ExtAll,
})
\end{lstlisting}

For standalone documents, \texttt{MakeTitlePage} and \texttt{MakeTOC} entries in the definition block automatically place \texttt{\textbackslash{}maketitle} and \texttt{\textbackslash{}tableofcontents} at the start of the body.

An indented example showing raw AST traversal:

\begin{verbatim}
doc := parser.Parse([]byte("# Hello"), 0)
visitor.Walk(doc, myVisitor)
\end{verbatim}

\subsection{Special Characters}

LaTeX special characters in Markdown source are automatically escaped. For example, a 50\% discount on the price of \$10.00 \& free shipping becomes \texttt{50\textbackslash{}\% discount on the price of \textbackslash{}\$10.00 \textbackslash{}\& free shipping} in the output. The \texttt{\_} in identifiers like \texttt{my\_variable} is escaped to \texttt{my\textbackslash{}\_variable}.


\noindent\rule{\linewidth}{0.4pt}

\section{Conclusion}

As demonstrated by Figures~\ref{fig:fig_pipeline} and~\ref{fig:fig_ast} and Tables~\ref{tab:tab_comparison} and~\ref{tab:tab_performance}, MarkTeX~\cite{gruber2004markdown} achieves competitive performance with a clean, extensible design. The five-step extension pattern~\cite{commonmark2021spec} ensures that new constructs can be added without modifying existing code.

Future work will address the CommonMark delimiter-stack algorithm~\cite{commonmark2021spec} for fully compliant emphasis parsing.


\noindent\rule{\linewidth}{0.4pt}

See~\cite{knuth1986texbook},~\cite{lamport1994latex},~\cite{gruber2004markdown}, and~\cite{commonmark2021spec} for full bibliographic details.


\end{document}
